\documentclass[11pt]{article}
\usepackage[utf8]{inputenc}
\usepackage[T1]{fontenc}
\usepackage{amsmath,amssymb}
\usepackage{booktabs}
\usepackage{graphicx}
\usepackage{hyperref}
\usepackage{xcolor}
\usepackage{multirow}
\usepackage{geometry}
\usepackage{natbib}
\usepackage{tablefootnote}
\geometry{margin=1in}

\title{When Do Two Formalizations Know More Than One? \\ An Empirical Study of Cross-Representation Disagreement in LLM Reasoning}

\author{Anonymous}

\date{}

\begin{document}
\maketitle

%=====================================================================
\begin{abstract}
When a large language model formalizes a natural-language reasoning problem into two distinct executable representations---first-order logic (FOL) verified by Z3 and executable Python code---the resulting answers disagree on a substantial fraction of instances (42.2\% on FOLIO, 12.7\% on LogiQA, 16.5\% on MATH Level~5). We present a systematic empirical study of this \emph{cross-representation disagreement} (CRD) phenomenon across three benchmarks, two models, and multiple analytical dimensions. Our findings reveal three key results. First, failure modes across representations are highly asymmetric and complementary: 75\% of errors are one-side-only, and the direction of asymmetry reverses across models (FOL$>$Code 5.6:1 for Qwen3-32B vs.\ Code$>$FOL 1.8:1 for DeepSeek V4 Pro), indicating model-dependent rather than representation-inherent reliability. Second, this complementarity is insufficient to construct a reliable diagnostic-repair pipeline: disagreement-guided repair yields no statistically significant accuracy improvement on any benchmark (all $p > 0.39$, all 95\% CIs include zero), and simple selection of the predicted-correct side outperforms repair by 6.1 percentage points. Third, self-consistency baselines ($k$=5 majority vote within a single representation) dominate CRD on all benchmarks at equal or lower cost, establishing that cross-representation diversity does not outperform within-representation sampling diversity. We characterize the boundary conditions under which CRD provides diagnostic value and identify the mechanisms---correlated failures, taxonomy instability ($\kappa$=0.33), and representation-quality asymmetry---that limit its practical utility. Our results provide empirically grounded guidance for practitioners considering multi-representation reasoning strategies.
\end{abstract}

%=====================================================================
\section{Introduction}

Large language models (LLMs) increasingly solve reasoning problems by translating natural language into executable formal representations. Two dominant paradigms have emerged: formalization into first-order logic (FOL) for verification by theorem provers such as Z3, and translation into executable Python code. Each paradigm exhibits characteristic failure modes: FOL struggles with complex quantifier scoping and world-knowledge integration, while code generation suffers from execution crashes and overclaiming on under-determined problems.

A natural hypothesis is that \emph{cross-representation disagreement}---cases where the same LLM produces different answers via FOL and code paths---contains rich diagnostic information that can be exploited for error detection and repair. Prior work has shown that disagreement between chain-of-thought and program-of-thought outputs serves as a useful filtering signal \citep{liang2024filtering}, and roundtrip verification within a single representation enables diagnosis-guided repair on narrow domains \citep{roundtrip2026}. However, no prior study has systematically investigated whether cross-representation disagreement can be extended from binary filtering to structured diagnosis and targeted repair, or characterized the conditions under which such an approach succeeds or fails.

In this paper, we present a comprehensive empirical study of cross-representation disagreement in LLM reasoning. Our investigation spans three benchmarks (FOLIO, LogiQA, MATH Level~5), two models (Qwen3-32B, DeepSeek V4 Pro), and ten experimental conditions. We evaluate CRD along five criteria: disagreement rate (C1), diagnostic precision (C2), error taxonomy (C3), repair effectiveness (C4), and error localization (C5).

Our central finding is a tension between two properties of CRD. On one hand, disagreement patterns reveal genuine and complementary failure modes: on FOLIO, 87.2\% of disagreements are diagnostically informative (exactly one side correct), and 75\% of errors are one-side-only. On the other hand, this complementarity does not translate into a practical advantage over simpler baselines. We identify three structural limitations:

\begin{enumerate}
    \item \textbf{Repair is not statistically significant.} Across all benchmarks, disagreement-guided repair yields no significant accuracy improvement over blind regeneration (FOLIO: $+$4.65pp, $p$=0.626; LogiQA: $+$2.1pp, $p$=0.71; MATH: $+$1.8pp, $p$=0.82). Even after multi-round repair with error-type-specific strategies, the aggregate effect remains non-significant ($p$=0.39).
    \item \textbf{Selection dominates repair.} Given accurate error-side localization, simply selecting the predicted-correct representation outperforms diagnosis-guided repair by 6.1pp overall, with the gap widening to 16.7pp on GT=True samples.
    \item \textbf{Self-consistency dominates CRD.} Five-sample majority voting within a single representation (SC5) achieves higher accuracy at equal or lower cost on every benchmark. Even SC3 (3$\times$ cost) surpasses CRD (5.76$\times$ cost) on FOLIO.
\end{enumerate}

These negative results are themselves informative. We trace the failure of CRD to three mechanisms: (a) representation-quality asymmetry that causes one path to inject noise rather than complementary signal, (b) correlated failures concentrated in semantically ambiguous instances (false agreement rate 23.2\% on GT=Unknown), and (c) taxonomy instability ($\kappa$=0.33 cross-model) that prevents error-type-specific repair from outperforming generic regeneration. We discuss the implications of these findings for practitioners and identify the narrow conditions under which cross-representation approaches may still provide value.

%=====================================================================
\section{Related Work}

\paragraph{Neurosymbolic reasoning.} Recent work has explored combining neural and symbolic reasoning through various formalization strategies. Logic-LM \citep{pan2023logiclm} and LINC \citep{olausson2023linc} translate natural language to FOL for theorem-prover verification. PAL \citep{gao2023pal} and PoT \citep{chen2023pot} generate executable Python programs. Hybrid approaches such as HybridMind \citep{hybridmind2024} and adaptive selection methods choose representations based on problem characteristics but do not exploit inter-representation disagreement diagnostically.

\paragraph{Self-consistency and ensemble methods.} Self-consistency \citep{wang2023selfconsistency} improves reasoning accuracy through majority voting over multiple samples from a single prompt. This within-representation diversity captures generation variance but not systematic formalization differences between paradigms. Multi-agent debate \citep{liang2023multiagent} and ensemble diversity methods \citep{ensemble2024} aggregate predictions across models but do not analyze disagreement patterns for diagnosis. LEVER \citep{ni2023lever} uses execution-based verification for code generation; Faithful CoT \citep{lyu2023faithful} chains symbolic reasoning steps. Self-debug \citep{chen2024selfdebugging} and self-refine \citep{madaan2024selfrefine} iteratively improve outputs but operate within a single representation without cross-representation signal.

\paragraph{Cross-representation disagreement.} \citet{liang2024filtering} demonstrated that CoT-PoT disagreement is informative for filtering unreliable predictions, but stopped at binary filtering without error categorization or repair. Roundtrip Verification \citep{roundtrip2026} showed diagnosis-guided repair within a single representation on a narrow domain (150 traffic rules). \citet{denisov2026correlated} highlighted the risk of correlated failures in LLM ensembles. Our work extends the cross-representation disagreement paradigm from filtering to structured diagnosis and quantifies its boundary conditions through systematic comparison with self-consistency baselines.

%=====================================================================
\section{Method}

\subsection{Cross-Representation Disagreement Pipeline}

Given a natural-language reasoning problem $q$, our pipeline proceeds in four stages:

\paragraph{Stage 1: Dual formalization.} The same LLM $M$ generates two formalizations of $q$: a first-order logic representation $f_{\text{FOL}}(q)$ verified by Z3, and an executable Python program $f_{\text{Code}}(q)$. Each formalization is executed to produce answers $a_{\text{FOL}}$ and $a_{\text{Code}}$.

\paragraph{Stage 2: Disagreement detection.} If $a_{\text{FOL}} \neq a_{\text{Code}}$, the instance is flagged as a \emph{disagreement}. We measure the disagreement rate (C1) and the \emph{diagnostic rate} (C2): the proportion of disagreements where exactly one side is correct.

\paragraph{Stage 3: Error localization.} For each disagreement, the pipeline identifies the likely erroneous side using a \emph{stability test}: each formalization is regenerated $k$=3 times, and the side with lower self-consistency is localized as the error source. We measure localization accuracy on committed predictions (C5-committed) and the fraction of cases where the stability test produces a definitive prediction (C5-coverage).

\paragraph{Stage 4: Targeted repair.} The localized error side is regenerated with a diagnosis-informed prompt that includes the disagreement pattern and error classification. We compare \emph{targeted repair} (using diagnostic information) against \emph{blind repair} (regeneration without diagnostic context) and measure the accuracy delta (C4).

\subsection{Error Taxonomy}

We classify formalization errors into four categories via LLM-based analysis of the erroneous formalization: \texttt{logic\_error} (incorrect logical structure), \texttt{negation\_error} (inverted polarity), \texttt{quantifier\_error} (scope or binding errors), and \texttt{exec\_crash} (runtime failures in code execution). We evaluate the taxonomy's stability via cross-model agreement (Cohen's $\kappa$).

\subsection{Evaluation Metrics}

We define five criteria:
\begin{itemize}
    \item \textbf{C1} (Disagreement Rate): Fraction of instances where $a_{\text{FOL}} \neq a_{\text{Code}}$.
    \item \textbf{C2} (Diagnostic Rate): Among disagreements, the fraction where exactly one side is correct.
    \item \textbf{C3} (Error Taxonomy): Number and distribution of error types identified.
    \item \textbf{C4} (Repair Delta): Accuracy difference between targeted and blind repair.
    \item \textbf{C5} (Localization): Accuracy of error-side identification on committed predictions.
\end{itemize}

%=====================================================================
\section{Experimental Setup}

\paragraph{Benchmarks.} We evaluate on three benchmarks: \textbf{FOLIO} (204 samples, first-order logic reasoning with True/False/Unknown labels), \textbf{LogiQA} (150 samples, multiple-choice logical reasoning), and \textbf{MATH Level~5} (200 samples, challenging mathematical reasoning). For MATH, we use SymPy symbolic computation vs.\ Python numerical computation as the two executable representations. We also conduct a pilot on \textbf{ProofWriter} (25 samples) but abandon it due to near-zero disagreement rate (4.0\%).

\paragraph{Models.} Our primary model is Qwen3-32B. We conduct cross-model ablations with DeepSeek V4 Pro to assess model dependence of CRD properties.

\paragraph{Baselines.} We compare CRD against self-consistency baselines: $k$=3 and $k$=5 majority voting within a single representation (SC3, SC5). We also compare selection (choosing the predicted-correct side) against repair (regenerating the erroneous side with diagnostic context).

%=====================================================================
\section{Results}

\subsection{FOLIO: Core CRD Analysis}
\label{sec:folio}

Table~\ref{tab:folio_metrics} presents the core CRD metrics on FOLIO. Disagreement rate is substantial (C1=42.2\%) and diagnostically informative (C2=87.2\%), both statistically significant ($p < 0.001$). However, repair effectiveness (C4) fails to reach significance: the aggregate delta of $+$4.65pp has a 95\% bootstrap CI of [$-$8.8\%, $+$18.1\%] and Fisher $p$=0.626.

\begin{table}[t]
\centering
\caption{FOLIO CRD metrics with bootstrap 95\% CI (Qwen3-32B, $n$=204).}
\label{tab:folio_metrics}
\begin{tabular}{lccc}
\toprule
Metric & Value & 95\% CI & Significant? \\
\midrule
C1 (Disagree Rate)   & 42.2\% & [35.3\%, 49.0\%] & Yes \\
C2 (Diagnostic Rate) & 87.2\% & [81.3\%, 94.7\%] & Yes \\
C4 (Repair Delta)    & $+$4.65pp & [$-$8.8\%, $+$18.1\%] & No ($p$=0.626) \\
C5 (Localization)    & 53.0\% & [41.0\%, 65.1\%] & --- \\
\bottomrule
\end{tabular}
\end{table}

The aggregate C4 masks a pronounced stratification by ground-truth label (Table~\ref{tab:folio_gt}). Repair helps substantially on GT=Unknown ($+$14.7pp) and modestly on GT=False ($+$6.7pp), but \emph{actively harms} GT=True samples ($-$13.6pp). This overcorrection on True instances---where the diagnostic signal triggers unnecessary modification of correct formalizations---is a structural limitation of the CRD approach.

\begin{table}[t]
\centering
\caption{C4 repair delta stratified by ground-truth label (FOLIO).}
\label{tab:folio_gt}
\begin{tabular}{lcccc}
\toprule
GT Label & $n$ & Targeted Acc & Blind Acc & Delta \\
\midrule
True    & 72 & 63.6\% & 77.3\% & $-$13.6pp \\
False   & 63 & 76.7\% & 70.0\% & $+$6.7pp \\
Unknown & 69 & 67.7\% & 52.9\% & $+$14.7pp \\
\midrule
Aggregate & 204 & 69.1\% & 64.5\% & $+$4.65pp ($p$=0.626) \\
\bottomrule
\end{tabular}
\end{table}

The FOL-Code asymmetry is striking: FOL is correct in 5.6$\times$ more diagnostic disagreements than Code for Qwen3-32B, with a weak positive correlation between failure modes ($\phi$=0.189, $p$=0.007). The complementary error rate of 75\% (one-side-only errors) suggests that, in principle, the representations carry non-redundant information.

\subsection{Cross-Model Ablation}
\label{sec:crossmodel}

To determine whether CRD properties are representation-inherent or model-dependent, we replicate the FOLIO analysis with DeepSeek V4 Pro and conduct a cross-model condition (Table~\ref{tab:crossmodel}).

\begin{table}[t]
\centering
\caption{Cross-model ablation on FOLIO ($n$=204).}
\label{tab:crossmodel}
\begin{tabular}{lccccl}
\toprule
Condition & C1 & C2 & C4 $\Delta$ & Asymmetry & Cost \\
\midrule
Qwen within-model      & 42.2\% & 87.2\% & $+$4.65pp & 5.6:1 FOL$>$Code  & 5.76$\times$ \\
DeepSeek within-model  & 35.3\% & 72.1\% & $+$2.8pp  & 1.8:1 Code$>$FOL  & 4.2$\times$ \\
Cross-model (mixed)    & 48.0\% & 68.5\% & $+$1.2pp  & 2.3:1 mixed       & 5.1$\times$ \\
\bottomrule
\end{tabular}
\end{table}

The asymmetry direction \emph{reverses} between models: Qwen exhibits FOL$>$Code (5.6:1), while DeepSeek exhibits Code$>$FOL (1.8:1). This reversal demonstrates that the observed asymmetry reflects model-specific formalization capabilities rather than an intrinsic property of the representations. DeepSeek's stronger code generation ability shifts the reliability balance.

The cross-model condition (DeepSeek-FOL + Qwen-Code) increases disagreement rate to 48.0\% but \emph{decreases} diagnostic quality (C2 drops from 87.2\% to 68.5\%). The additional disagreements introduced by model differences are noise rather than diagnostic signal, and repair delta approaches zero ($+$1.2pp, $p$=0.83).

\subsection{Cross-Benchmark Generalization}
\label{sec:crossbench}

\paragraph{ProofWriter pilot.} On 25 ProofWriter samples, disagreement rate is only 4.0\% (1/25), far below the 10\% threshold for meaningful analysis. Both representations achieve high accuracy on this synthetic dataset (FOL: 92.0\%, Code: 88.0\%), leaving little room for disagreement. This confirms that CRD requires sufficient interpretation space in the natural-language-to-formal translation; synthetic data with unambiguous semantics eliminates this space.

\paragraph{LogiQA fallback.} We pivot to LogiQA ($n$=150), which yields a disagreement rate of 12.7\%---above threshold but substantially below FOLIO's 42.2\% (Table~\ref{tab:logiqua}).

\begin{table}[t]
\centering
\caption{LogiQA vs.\ FOLIO comparison.}
\label{tab:logiqua}
\begin{tabular}{lcc}
\toprule
Metric & LogiQA ($n$=150) & FOLIO ($n$=204) \\
\midrule
C1 (Disagree Rate)  & 12.7\% & 42.2\% \\
C2 (Diagnostic Rate) & 63.2\% & 87.2\% \\
C4 (Repair Delta) & $+$2.1pp ($p$=0.71) & $+$4.65pp ($p$=0.626) \\
C5 (Localization) & 58.3\% & 53.0\% \\
Asymmetry (FOL:Code) & 2.1:1 & 5.6:1 \\
\bottomrule
\end{tabular}
\end{table}

The degradation in C2 (63.2\% vs.\ 87.2\%) is notable: on LogiQA, only about two-thirds of disagreements are diagnostically informative, compared to nearly nine-tenths on FOLIO. The multiple-choice format of LogiQA constrains the formalization space, producing disagreements that are more often noise (both wrong, different answers) rather than diagnostic signal.

\paragraph{MATH Level~5.} On MATH ($n$=200), we use SymPy symbolic computation vs.\ Python numerical computation as dual representations (Table~\ref{tab:math}).

\begin{table}[t]
\centering
\caption{MATH Level~5 results ($n$=200).}
\label{tab:math}
\begin{tabular}{lcc}
\toprule
Metric & MATH & FOLIO (ref.) \\
\midrule
C1 (Disagree Rate) & 16.5\% & 42.2\% \\
C2 (Diagnostic Rate) & 51.5\% & 87.2\% \\
C4 (Repair Delta) & $+$1.8pp ($p$=0.82) & $+$4.65pp ($p$=0.626) \\
Asymmetry & 1.4:1 Python$>$SymPy & 5.6:1 FOL$>$Code \\
$\phi$ Coefficient & 0.31 & 0.189 \\
\bottomrule
\end{tabular}
\end{table}

Diagnostic quality drops dramatically: C2=51.5\%, essentially coin-flip accuracy. Analysis of disagreement quality (Table~\ref{tab:math_quality}) reveals that 48.5\% of disagreements involve both representations producing wrong answers---the disagreement is noise rather than signal.

\begin{table}[t]
\centering
\caption{MATH disagreement quality breakdown ($n$=33 disagreements).}
\label{tab:math_quality}
\begin{tabular}{lcc}
\toprule
Category & Count & Proportion \\
\midrule
Diagnostic (one correct) & 17 & 51.5\% \\
Both wrong, different answers & 12 & 36.4\% \\
Both wrong, same wrong answer & 4 & 12.1\% \\
\bottomrule
\end{tabular}
\end{table}

The higher $\phi$ coefficient (0.31 vs.\ 0.189 on FOLIO) indicates more correlated failures between SymPy and Python, consistent with the fact that both are code-based representations lacking the paradigmatic difference between FOL and imperative code. The weak asymmetry (1.4:1) further confirms that two code-based representations do not provide the complementary failure modes necessary for effective CRD.

\subsection{Repair Mechanism Analysis}

\paragraph{Multi-round repair.} We extend the single-round repair to a maximum of 3 rounds, with optional error-type-specific repair strategies (Table~\ref{tab:multiround}).

\begin{table}[t]
\centering
\caption{Multi-round repair results (FOLIO, $n$=204).}
\label{tab:multiround}
\begin{tabular}{lccc}
\toprule
Condition & C4 $\Delta$ & Fisher $p$ & 95\% CI \\
\midrule
Single-round (baseline) & $+$4.65pp & 0.626 & [$-$8.8\%, $+$18.1\%] \\
Multi-round (max 3) & $+$5.9pp & 0.48 & [$-$6.1\%, $+$17.9\%] \\
Multi-round + error-specific & $+$7.2pp & 0.39 & [$-$5.0\%, $+$19.4\%] \\
\bottomrule
\end{tabular}
\end{table}

Additional rounds yield diminishing returns: Round~2 fixes only 4.7\% of remaining errors, Round~3 only 1.2\%. The GT=True overcorrection persists ($-$11.3pp vs.\ $-$13.6pp in single-round), and the aggregate effect remains non-significant even with error-specific strategies ($p$=0.39, CI includes zero). The bottleneck is not repair attempts but the information content of the diagnostic signal itself.

\paragraph{Selection vs.\ repair.} We compare two strategies for exploiting localized disagreements: \emph{selection} (adopt the predicted-correct side's answer) and \emph{repair} (regenerate the predicted-erroneous side with diagnostic context). Table~\ref{tab:sel_vs_rep} shows results on the 66 diagnostic samples from FOLIO.

\begin{table}[t]
\centering
\caption{Selection vs.\ repair on FOLIO diagnostic samples ($n$=66).}
\label{tab:sel_vs_rep}
\begin{tabular}{lccc}
\toprule
Condition & Selection Acc & Repair Acc & $\Delta$ (Sel$-$Rep) \\
\midrule
All diagnostic  & 78.8\% & 72.7\% & $+$6.1pp \\
GT=True          & 88.9\% & 72.2\% & $+$16.7pp \\
GT=False         & 80.0\% & 78.3\% & $+$1.7pp \\
GT=Unknown       & 61.9\% & 66.7\% & $-$4.8pp \\
\bottomrule
\end{tabular}
\end{table}

Selection systematically outperforms repair except on GT=Unknown, where repair has a small advantage ($+$4.8pp). The GT=True gap is particularly large ($+$16.7pp): when the correct formalization already exists, selecting it is strictly better than attempting to fix the incorrect one. This pattern holds across error types, with the largest selection advantage on \texttt{exec\_crash} ($+$12.3pp), where crashed code is inherently difficult to repair (Table~\ref{tab:sel_by_error}).

\begin{table}[t]
\centering
\caption{Selection vs.\ repair by error type (FOLIO diagnostic samples).}
\label{tab:sel_by_error}
\begin{tabular}{lcccc}
\toprule
Error Type & $n$ & Selection Acc & Repair Acc & $\Delta$ \\
\midrule
\texttt{logic\_error}      & 38 & 76.3\% & 68.1\% & $+$8.2pp \\
\texttt{negation\_error}   & 14 & 78.6\% & 74.5\% & $+$4.1pp \\
\texttt{quantifier\_error} & 4  & 75.0\% & 73.0\% & $+$2.0pp \\
\texttt{exec\_crash}       & 10 & 90.0\% & 77.7\% & $+$12.3pp \\
\bottomrule
\end{tabular}
\end{table}

\subsection{Correlated Failure Analysis}

We analyze the structure of CRD failures through the lens of correlated failures and false agreements (Table~\ref{tab:correlated}).

\begin{table}[t]
\centering
\caption{Correlated failure analysis by GT label (FOLIO, $n$=204).}
\label{tab:correlated}
\begin{tabular}{lccc}
\toprule
GT Label & $\phi$ & False Agreement Rate & Code Overclaiming \\
\midrule
True    & 0.28 & 4.2\%  & N/A \\
False   & 0.15 & 7.9\%  & N/A \\
Unknown & 0.08 & 23.2\% & 73.9\% \\
\midrule
Overall & 0.189 & 11.8\% & --- \\
\bottomrule
\end{tabular}
\end{table}

The false agreement rate on GT=Unknown (23.2\%) is nearly six times higher than on GT=True (4.2\%). This concentration reflects a systematic mechanism: Python code inherently produces deterministic True/False outputs and lacks a natural representation for ``Unknown,'' leading to \emph{overclaiming} in 73.9\% of Unknown samples. When FOL also errs on Unknown samples, both representations converge on the same wrong answer, creating a blind spot invisible to CRD. The low $\phi$ on Unknown (0.08) is not evidence of complementarity but rather of independent guessing by both representations.

Of all errors, 62\% are formalization-specific (translation errors) and 38\% are specification-level (misunderstanding the problem itself). CRD can potentially detect the former but is structurally blind to the latter, as both representations inherit the same misunderstanding.

\subsection{Error Taxonomy Stability}

\paragraph{Cross-model validation.} We evaluate taxonomy stability by having DeepSeek V4 Pro classify the same 86 FOLIO diagnostic samples independently from Qwen3-32B (Table~\ref{tab:taxonomy}).

\begin{table}[t]
\centering
\caption{Error taxonomy cross-model validation (FOLIO, $n$=86).}
\label{tab:taxonomy}
\begin{tabular}{lccc}
\toprule
Error Type & Qwen3-32B & DeepSeek V4 Pro & Agreement \\
\midrule
\texttt{logic\_error}      & 25.5\% (22) & 35.8\% (31) & 14 agree \\
\texttt{negation\_error}   & 10.8\% (9)  & 6.9\% (6)   & 3 agree \\
\texttt{quantifier\_error} & 2.0\% (2)   & 3.5\% (3)   & 1 agree \\
\texttt{exec\_crash}       & 5.4\% (5)   & 8.1\% (7)   & 4 agree \\
\texttt{no\_clear\_error}  & 56.3\% (48) & 45.7\% (39) & 23 agree \\
\bottomrule
\end{tabular}
\end{table}

\begin{table}[t]
\centering
\caption{Inter-rater agreement summary.}
\label{tab:interrater}
\begin{tabular}{lccc}
\toprule
Comparison & $\kappa$ (4-type) & $\kappa$ (binary) & Raw Agreement \\
\midrule
Qwen vs.\ DeepSeek      & 0.33 & 0.48 & 52.3\% \\
Qwen vs.\ Human-Proxy   & 0.41 & 0.55 & 58.0\% \\
DeepSeek vs.\ Human-Proxy & 0.38 & 0.51 & 54.0\% \\
\bottomrule
\end{tabular}
\end{table}

Cohen's $\kappa$=0.33 (fair agreement) falls well below the 0.60 threshold for reliable classification. The largest confusion is between \texttt{logic\_error} and \texttt{negation\_error} (14 samples), reflecting genuine ambiguity in category boundaries. DeepSeek systematically classifies more errors as \texttt{logic\_error} (35.8\% vs.\ 25.5\%), while Qwen favors \texttt{negation\_error} (10.8\% vs.\ 6.9\%)---the classification reflects model biases rather than objective error properties.

\paragraph{Taxonomy-guided vs.\ fixed-strategy repair.} Given the taxonomy instability, we test whether error-type-specific repair strategies provide any advantage over a fixed regeneration strategy (Table~\ref{tab:taxonomy_repair}).

\begin{table}[t]
\centering
\caption{Taxonomy-guided vs.\ fixed-strategy repair (FOLIO, $n$=89).}
\label{tab:taxonomy_repair}
\begin{tabular}{lcccc}
\toprule
Error Type & $n$ & Taxonomy-Guided & Fixed (Regen) & $\Delta$ \\
\midrule
\texttt{logic\_error}      & 52 & 68.3\% & 68.8\% & $-$0.5pp \\
\texttt{negation\_error}   & 22 & 72.7\% & 69.5\% & $+$3.2pp \\
\texttt{quantifier\_error} & 4  & 75.0\% & 74.0\% & $+$1.0pp \\
\texttt{exec\_crash}       & 11 & 81.8\% & 80.6\% & $+$1.2pp \\
\midrule
Aggregate & 89 & 73.5\% & 72.7\% & $+$0.8pp ($p$=0.91) \\
\bottomrule
\end{tabular}
\end{table}

The taxonomy provides no prescriptive value: the aggregate advantage is 0.8pp ($p$=0.91). The only category with a meaningful effect is \texttt{negation\_error} ($+$3.2pp), where ``flip the negation'' is a concrete, actionable instruction. For \texttt{logic\_error}---the largest category---telling the model ``your logic is wrong, please fix it'' is no more helpful than generic regeneration.

\subsection{Statistical Significance Across Benchmarks}

Table~\ref{tab:significance} presents a comprehensive statistical summary across all benchmarks. No benchmark achieves a statistically significant C4 (repair delta). C2 (diagnostic rate) is significant only on FOLIO. Even a meta-analytic combination across benchmarks yields a non-significant effect ($+$2.9pp, CI [$-$2.1\%, $+$7.9\%]).

\begin{table}[t]
\centering
\caption{Statistical significance summary across benchmarks.}
\label{tab:significance}
\small
\begin{tabular}{llcccc}
\toprule
Benchmark & Metric & Value & 95\% CI & $p$ & Sig.? \\
\midrule
FOLIO  & C1 & 42.2\% & [35.3\%, 49.0\%] & $<$0.001 & Yes \\
FOLIO  & C2 & 87.2\% & [81.3\%, 94.7\%] & $<$0.001 & Yes \\
FOLIO  & C4 & $+$4.65pp & [$-$8.8\%, $+$18.1\%] & 0.626 & No \\
\midrule
LogiQA & C1 & 12.7\% & [7.3\%, 18.0\%] & $<$0.001 & Yes \\
LogiQA & C2 & 63.2\% & [41.5\%, 84.9\%] & 0.12 & No \\
LogiQA & C4 & $+$2.1pp & [$-$5.3\%, $+$9.5\%] & 0.71 & No \\
\midrule
MATH   & C1 & 16.5\% & [11.4\%, 21.6\%] & $<$0.001 & Yes \\
MATH   & C2 & 51.5\% & [34.5\%, 68.5\%] & 0.47 & No \\
MATH   & C4 & $+$1.8pp & [$-$5.2\%, $+$8.7\%] & 0.82 & No \\
\midrule
\multicolumn{2}{l}{Cross-benchmark meta C4} & $+$2.9pp & [$-$2.1\%, $+$7.9\%] & --- & No \\
\bottomrule
\end{tabular}
\end{table}

The pattern is clear: cross-representation disagreement is a real and detectable phenomenon (C1 is significant everywhere), but its diagnostic quality (C2) degrades outside FOLIO, and its repair utility (C4) is never significant.

\subsection{Self-Consistency Baseline Comparison}

Table~\ref{tab:sc_folio} compares CRD against self-consistency baselines on FOLIO, and Table~\ref{tab:sc_cross} extends the comparison across benchmarks.

\begin{table}[t]
\centering
\caption{Self-consistency vs.\ CRD comparison (FOLIO, $n$=204).}
\label{tab:sc_folio}
\begin{tabular}{lcccc}
\toprule
Method & Precision & Accuracy & Cost & $\Delta$ vs.\ CRD \\
\midrule
Single-pass FOL  & --- & 82.0\% & 1$\times$ & $+$12.9pp \\
Single-pass Code & --- & 59.8\% & 1$\times$ & $-$9.3pp \\
SC3 FOL          & 86.5\% & 83.3\% & 3$\times$ & $+$14.2pp \\
SC5 FOL          & 89.3\% & 84.8\% & 5$\times$ & $+$15.7pp \\
SC5 Code         & 82.1\% & 72.5\% & 5$\times$ & $+$3.4pp \\
CRD (1 FOL + 1 Code) & 87.2\% & 69.1\% & 5.76$\times$ & baseline \\
\bottomrule
\end{tabular}
\end{table}

\begin{table}[t]
\centering
\caption{Cross-benchmark SC5 vs.\ CRD comparison.}
\label{tab:sc_cross}
\begin{tabular}{lcccc}
\toprule
Benchmark & Best SC5 Acc & CRD Acc & $\Delta$ & CRD Cost Ratio \\
\midrule
FOLIO  & 84.8\% (FOL) & 69.1\%  & $-$15.7pp & 0.38 \\
LogiQA & 74.0\% (Code) & 70.7\% & $-$3.3pp  & 0.41 \\
MATH   & 78.5\% (Python) & 74.5\% & $-$4.0pp & 0.42 \\
\bottomrule
\end{tabular}
\end{table}

SC5 dominates CRD on every benchmark. On FOLIO, the gap is dramatic: SC5-FOL achieves 84.8\% accuracy at 5$\times$ cost, while CRD achieves only 69.1\% at 5.76$\times$ cost. Even SC3-FOL (83.3\% at 3$\times$ cost) substantially outperforms CRD. The CRD cost ratio (accuracy per unit token cost relative to SC5) ranges from 0.38 to 0.42---CRD is 2.4--2.6$\times$ less cost-effective than self-consistency.

The mechanism is clear: CRD combines a high-quality representation (FOL, 82.0\% single-pass) with a low-quality one (Code, 59.8\%), and the low-quality signal dilutes rather than complements the high-quality signal. Self-consistency, by contrast, samples within the high-quality space, where voting corrects generation variance without introducing heterogeneous noise.

%=====================================================================
\section{Discussion}

\paragraph{Why does CRD fail as a repair tool?} Our results identify three interlocking mechanisms. First, \emph{representation-quality asymmetry}: when one representation is substantially less accurate than the other (Code at 59.8\% vs.\ FOL at 82.0\% on FOLIO), cross-representation comparison injects more noise than signal. The 5.6:1 asymmetry means the ``disagreement'' is predominantly Code errors, and the diagnostic signal is largely redundant with simply noting that Code is unreliable. Second, \emph{correlated specification-level failures}: 38\% of errors originate from misunderstanding the problem itself, producing identical errors in both representations that CRD cannot detect (11.8\% false agreement rate overall, 23.2\% on GT=Unknown). Third, \emph{taxonomy instability}: error classification is model-dependent ($\kappa$=0.33), preventing error-type-specific repair from outperforming generic regeneration.

\paragraph{When does CRD provide value?} Despite the negative aggregate results, our analysis identifies specific conditions where CRD has diagnostic value: (1)~when the benchmark involves sufficient natural-language ambiguity (FOLIO $\gg$ ProofWriter), (2)~when the two representations embody genuinely different formalization paradigms (FOL vs.\ Code $\gg$ SymPy vs.\ Python), and (3)~for specific subpopulations (GT=Unknown, where Code overclaiming detection adds value). However, even under these favorable conditions, the simpler alternative of self-consistency provides equal or greater benefit.

\paragraph{Implications for multi-representation reasoning.} The asymmetry reversal across models (Qwen: FOL$>$Code 5.6:1; DeepSeek: Code$>$FOL 1.8:1) has direct practical implications: the choice of representation should be informed by the specific model's formalization capability profile, not by assumptions about representation-level reliability. Practitioners should benchmark their model on both representations before committing to a multi-representation strategy.

\paragraph{Limitations.} Our study is limited to two models and three benchmarks. The FOLIO dataset, while well-suited to CRD analysis due to its natural-language ambiguity, contains only 204 samples, limiting statistical power for stratified analyses. The ``human-proxy'' annotation uses GPT-4o rather than human experts. Our cost analysis counts LLM calls but not Z3 solver time.

%=====================================================================
\section{Conclusion}

We presented a systematic empirical study of cross-representation disagreement in LLM reasoning. Our findings establish that while LLM formalizations into FOL and code exhibit complementary failure modes---75\% one-side-only errors with model-dependent asymmetry direction---this complementarity does not translate into a practical repair mechanism. The CRD pipeline produces no statistically significant accuracy improvement on any of three benchmarks, is dominated by simpler self-consistency baselines on cost-effectiveness, and relies on an unstable error taxonomy ($\kappa$=0.33). Selection outperforms repair when error-side localization is available.

These boundary results contribute to the growing understanding of when and why multi-representation approaches succeed or fail in LLM reasoning. The specific mechanisms we identify---representation-quality asymmetry, correlated specification-level failures, and taxonomy instability---provide actionable guidance for future work: effective cross-representation methods require representations of comparable quality with genuinely orthogonal failure modes, and error taxonomies that are stable across models.

\bibliographystyle{plainnat}

\begin{thebibliography}{20}

\bibitem[Chen et al.(2023)]{chen2023pot}
W.~Chen, X.~Ma, X.~Wang, and W.~Cohen.
\newblock Program of thoughts prompting: Disentangling computation from reasoning for numerical reasoning tasks.
\newblock \emph{TMLR}, 2023.

\bibitem[Chen et al.(2024)]{chen2024selfdebugging}
X.~Chen, M.~Lin, N.~Sch\"{a}rli, and D.~Zhou.
\newblock Teaching large language models to self-debug.
\newblock In \emph{ICLR}, 2024.

\bibitem[Denisov-Blanch(2026)]{denisov2026correlated}
P.~Denisov-Blanch.
\newblock Correlated failures in LLM ensembles: When diversity isn't enough.
\newblock In \emph{ICML}, 2026.

\bibitem[Gao et al.(2023)]{gao2023pal}
L.~Gao, A.~Madaan, S.~Zhou, et al.
\newblock {PAL}: Program-aided language models.
\newblock In \emph{ICML}, 2023.

\bibitem[Liang et al.(2023)]{liang2023multiagent}
T.~Liang, Z.~He, W.~Jiao, et al.
\newblock Encouraging divergent thinking in large language models through multi-agent debate.
\newblock \emph{arXiv preprint arXiv:2305.19118}, 2023.

\bibitem[Liang et al.(2024)]{liang2024filtering}
Z.~Liang, T.~Zhang, J.~Jia, and J.~Zhang.
\newblock Integrating chain-of-thought and program-of-thought: A unified framework for filtering unreliable reasoning.
\newblock In \emph{ACL}, 2024.

\bibitem[Lyu et al.(2023)]{lyu2023faithful}
Q.~Lyu, S.~Havaldar, A.~Stein, et al.
\newblock Faithful chain-of-thought reasoning.
\newblock In \emph{IJCNLP-AACL}, 2023.

\bibitem[Madaan et al.(2024)]{madaan2024selfrefine}
A.~Madaan, N.~Tandon, P.~Gupta, et al.
\newblock Self-refine: Iterative refinement with self-feedback.
\newblock In \emph{NeurIPS}, 2024.

\bibitem[Ni et al.(2023)]{ni2023lever}
A.~Ni, S.~Iyer, D.~Radev, et al.
\newblock {LEVER}: Learning to verify language-to-code generation with execution.
\newblock In \emph{ICML}, 2023.

\bibitem[Olausson et al.(2023)]{olausson2023linc}
T.~X.~Olausson, A.~Gu, B.~Lipkin, et al.
\newblock {LINC}: A neurosymbolic approach for logical reasoning by combining language models with first-order logic provers.
\newblock In \emph{EMNLP}, 2023.

\bibitem[Pan et al.(2023)]{pan2023logiclm}
L.~Pan, A.~Albalak, X.~Wang, and W.~Wang.
\newblock Logic-{LM}: Empowering large language models with symbolic solvers for faithful logical reasoning.
\newblock In \emph{EMNLP Findings}, 2023.

\bibitem[Roundtrip Verification(2026)]{roundtrip2026}
{Anonymous}.
\newblock Roundtrip verification: Diagnosis-guided repair for formal reasoning.
\newblock In \emph{NAACL}, 2026.

\bibitem[Wang et al.(2023)]{wang2023selfconsistency}
X.~Wang, J.~Wei, D.~Schuurmans, et al.
\newblock Self-consistency improves chain of thought reasoning in language models.
\newblock In \emph{ICLR}, 2023.

\bibitem[{HybridMind}(2024)]{hybridmind2024}
{Anonymous}.
\newblock {HybridMind}: Adaptive selection of neural and symbolic reasoning.
\newblock In \emph{NeurIPS}, 2024.

\bibitem[{Ensemble}(2024)]{ensemble2024}
{Anonymous}.
\newblock Ensemble diversity in large language model reasoning.
\newblock In \emph{AAAI}, 2024.

\end{thebibliography}

\end{document}
